\small
\textbf{Abstract.}
Retail forecasting models that see the price and promotion plan can explain
each forecast as a baseline, what would sell with nothing planned, plus the
lift each part of the plan adds. Merchants read that split as fact and plan on
it. We set out eight checks on such explanations, fixed in advance: properties
that hold by construction, sensitivity to randomised weights, stability across
seeds, price direction, agreement with what sold, agreement with LightGBM, a
plan read-back, and agreement with a measured cause. In generated store data
with known effects, a world model whose explanation is built in passes C1 and
C3 to C6, and its lifts match the truth
(rank agreement \kaClWmTruth) when plans are made at random. When merchants plan
into high-demand weeks, it still passes them, forecasts more accurately,
and overstates the total lift by a third (\kaCoWmBiasLo{} to \kaCoWmBiasHi{}
times) and the display's by \kaCoWmDispLo{} to \kaCoWmDispHi{} times; LightGBM
overstates them too, and the two agree. Asked about a lasting change in the
regular price, the world model answers with the temporary-cut response,
$\kaClWmCsevenLo$ to $\kaClWmCsevenHi$ against a true $-0.5$, and only the
measured cause finds it. The price-direction check catches every model blind to
price but also flags LightGBM; an exploratory lever-by-lever comparison catches
lifts given to the wrong lever; the randomisation check separates nothing. On
Dominick's we find \zEpisodes{} spells in which a store's regular prices sat
apart from its zone peers, most between April 1993 and April 1994. Over the
\zScored{} scored spells, prices \zSizeLo\% to \zSizeHi\% off their peers, the
pooled effect on category units is $\zSlope$ over eight weeks (95\% interval
$\zSlopeLo$ to $\zSlopeHi$) and $\zFullSlope$ over the full spell ($\zFullLo$ to
$\zFullHi$), in line with the elasticity near $-0.3$ the chain's experiments
implied. These pooled intervals were added after a placebo at false start dates
gave $\zPlaceboSlope$ ($\zPlaceboLo$ to $\zPlaceboHi$) and showed single
episodes to be noise.
We then audited all \rVersions{} saved world-model and lift-model versions of
our companion paper on M5, Dominick's and Breakfast at the Frat. No version
passed every check on M5 or Dominick's. The baseline and lift add up to the
forecast in every version, but the lift was not zero when the audit planned
nothing on most Dominick's and Frat versions, for reasons we could not fully
trace; on Dominick's and M5 a deeper price cut
lowered the forecast in about half the item-weeks or more; and on M5 the lifts
agreed poorly between seeds although the forecasts agreed. Both models state
about twice the measured response to the zone price deviations, though the
measured interval is wide and overlaps both. Across versions,
the number of checks passed was not clearly related to accuracy (pooled
Spearman \rRhoPassedPool, $p = \rPPassed$, on the unit fixed in advance); the
check most strongly related to accuracy (C5) compares the forecast with sales. These results
hold for these models and data sets.

\vspace{4pt}
\textbf{Keywords:} retail forecasting; explanation; faithfulness; promotions;
price response; difference in differences; world models.
